% source-sha256: f668831a53e944be583111b9e6d49298e80bbc03ce4b11e64ad863c5cef40397
% Options for packages loaded elsewhere
\PassOptionsToPackage{unicode}{hyperref}
\PassOptionsToPackage{hyphens}{url}
\documentclass[
]{article}
\usepackage{xcolor}
\usepackage[letterpaper,margin=1in]{geometry}
\usepackage{amsmath,amssymb}
\setcounter{secnumdepth}{-\maxdimen} % remove section numbering
\usepackage{iftex}
\ifPDFTeX
  \usepackage[T1]{fontenc}
  \usepackage[utf8]{inputenc}
  \usepackage{textcomp} % provide euro and other symbols
\else % if luatex or xetex
  \usepackage{unicode-math} % this also loads fontspec
  \defaultfontfeatures{Scale=MatchLowercase}
  \defaultfontfeatures[\rmfamily]{Ligatures=TeX,Scale=1}
\fi
\usepackage{lmodern}
\ifPDFTeX\else
  % xetex/luatex font selection
\fi
% Use upquote if available, for straight quotes in verbatim environments
\IfFileExists{upquote.sty}{\usepackage{upquote}}{}
\IfFileExists{microtype.sty}{% use microtype if available
  \usepackage[]{microtype}
  \UseMicrotypeSet[protrusion]{basicmath} % disable protrusion for tt fonts
}{}
\makeatletter
\@ifundefined{KOMAClassName}{% if non-KOMA class
  \IfFileExists{parskip.sty}{%
    \usepackage{parskip}
  }{% else
    \setlength{\parindent}{0pt}
    \setlength{\parskip}{6pt plus 2pt minus 1pt}}
}{% if KOMA class
  \KOMAoptions{parskip=half}}
\makeatother
\usepackage{longtable,booktabs,array}
\usepackage{caption}
\captionsetup[table]{skip=6pt}
\newcounter{none} % for unnumbered tables
\usepackage{calc} % for calculating minipage widths
% Correct order of tables after \paragraph or \subparagraph
\usepackage{etoolbox}
\makeatletter
\patchcmd\longtable{\par}{\if@noskipsec\mbox{}\fi\par}{}{}
\makeatother
% Allow footnotes in longtable head/foot
\IfFileExists{footnotehyper.sty}{\usepackage{footnotehyper}}{\usepackage{footnote}}
\makesavenoteenv{longtable}
\setlength{\emergencystretch}{3em} % prevent overfull lines
\providecommand{\tightlist}{%
  \setlength{\itemsep}{0pt}\setlength{\parskip}{0pt}}
% Release-document layout safeguards injected by scripts/generate_docs.sh.
\usepackage{fvextra}
\fvset{breaklines=true,breakanywhere=true,fontsize=\small}
\RecustomVerbatimEnvironment{verbatim}{Verbatim}{breaklines=true,breakanywhere=true,fontsize=\small}
\setlength{\emergencystretch}{3em}
\AtBeginDocument{\sloppy}
% Workflow-figure support: mermaid blocks in the canonical Markdown are mapped
% to hand-authored TikZ figures by scripts/pandoc_bounded_tables.lua.
\usepackage{tikz}
\usetikzlibrary{shapes.geometric,arrows.meta,positioning,fit,backgrounds,calc}
\definecolor{bdiStep}{HTML}{EAF2FB}
\definecolor{bdiEdge}{HTML}{365F91}
\definecolor{bdiText}{HTML}{17253A}
\definecolor{bdiGate}{HTML}{FFF5D6}
\definecolor{bdiGateEdge}{HTML}{9B6B00}
\definecolor{bdiOpt}{HTML}{F0EAF7}
\definecolor{bdiOptEdge}{HTML}{5B3E86}
\definecolor{bdiOut}{HTML}{E6F4EA}
\definecolor{bdiOutEdge}{HTML}{1E6B3A}
\tikzset{
  bdi/.style={draw=bdiEdge,line width=1pt,fill=bdiStep,text=bdiText,rounded corners=3pt,
              align=center,inner sep=4pt,font=\scriptsize\sffamily},
  bdigate/.style={bdi,draw=bdiGateEdge,fill=bdiGate},
  bdiopt/.style={bdi,draw=bdiOptEdge,fill=bdiOpt,dashed},
  bdiout/.style={bdi,draw=bdiOutEdge,fill=bdiOut},
  bdiflow/.style={-{Stealth[length=2mm]},draw=bdiEdge,line width=0.9pt},
  bdiback/.style={bdiflow,dashed},
  bdilbl/.style={font=\tiny\sffamily,text=bdiText,inner sep=1.5pt,fill=white},
}
\usepackage{bookmark}
\IfFileExists{xurl.sty}{\usepackage{xurl}}{} % add URL line breaks if available
\urlstyle{same}
% fallback for those not using the hyperref driver hyperxmp:
\makeatletter
\@ifundefined{xmpquote}{\newcommand{\xmpquote}[1]{#1}}{}
\makeatother
\hypersetup{
  pdftitle={Business Document Ingestion},
  pdfauthor={Christopher Melhauser and theonlymuffinbot},
  hidelinks,
  pdfcreator={LaTeX via pandoc}}

\title{Business Document Ingestion}
\usepackage{etoolbox}
\makeatletter
\providecommand{\subtitle}[1]{% add subtitle to \maketitle
  \apptocmd{\@title}{\par {\large #1 \par}}{}{}
}
\makeatother
\subtitle{Client Output, Reporting, and CRM Access Overview}
\author{Christopher Melhauser and theonlymuffinbot}
\date{September 3, 2026}

\begin{document}
\maketitle

\section{Client Output, Reporting, and CRM Access
Overview}\label{client-output-reporting-and-crm-access-overview}

\subsection{Read this first}\label{read-this-first}

This guide explains what a completed document-processing engagement can
give you, what each output is for, and how approved data can later be
searched, reported on, exported, or prepared for a CRM. It is written
for business users; you do not need to know the internal processing
steps to use the results.

This edition also covers the optional visual-intake pilot and
source-only handoff implemented in this checkout. It does not claim a
deployed client service or a completed merge. The stable release remains
1.0.0; current repository and verification state is recorded separately
in \texttt{HANDOFF.md}.

The important distinction is simple:

{\def\LTcaptype{none} % do not increment counter
\begin{longtable}[]{@{}
  >{\raggedright\arraybackslash}p{(\linewidth - 4\tabcolsep) * \real{0.3333}}
  >{\raggedright\arraybackslash}p{(\linewidth - 4\tabcolsep) * \real{0.3333}}
  >{\raggedright\arraybackslash}p{(\linewidth - 4\tabcolsep) * \real{0.3333}}@{}}
\toprule\noalign{}
\begin{minipage}[b]{\linewidth}\raggedright
If the final review is\ldots{}
\end{minipage} & \begin{minipage}[b]{\linewidth}\raggedright
You receive\ldots{}
\end{minipage} & \begin{minipage}[b]{\linewidth}\raggedright
What it means
\end{minipage} \\
\midrule\noalign{}
\endhead
\bottomrule\noalign{}
\endlastfoot
\textbf{Open or blocked} & A review package and an exception register &
Questions still need an answer. The data is useful for review, but it is
not released as approved operational facts. \\
\textbf{Clear} & An approved data package, reporting/retrieval snapshot,
and optional CRM-ready files & Every included fact is provenance-linked,
review-clear, and tied to the exact source snapshot. \\
\end{longtable}
}

The system does not hide an unresolved item to make a spreadsheet look
clean. It records the question, reason, source, and next step. A
completed pipeline is therefore more than a collection of extracted
rows: it is a traceable answer to what was found, what was proven, and
what still needs a decision.

\subsection{The end product at a
glance}\label{the-end-product-at-a-glance}

A client delivery is organized into five practical layers. Depending on
the engagement and final gate, some layers may be delivered as blocked
findings rather than approved data. That distinction is visible in the
package.

\begin{enumerate}
\def\labelenumi{\arabic{enumi}.}
\tightlist
\item
  \textbf{Source and evidence layer.} Original documents, page
  references, hashes, and citations that let a user trace an important
  value back to its source.
\item
  \textbf{Review and control layer.} The review queue, validation,
  arithmetic, attribution, completeness, and exception results. This
  explains what was checked and what was not proven.
\item
  \textbf{Approved data layer.} A canonical JSON export containing only
  review-clear, provenance-linked facts, plus an exclusion artifact for
  every withheld record or field.
\item
  \textbf{Reporting and retrieval layer.} A read-only, immutable data
  snapshot that supports governed search, record lookup, account/contact
  cards, analysis, standard reports, and bounded exports.
\item
  \textbf{CRM handoff layer.} Checksummed no-send staging files and
  optional common CRM-import files. These prepare a controlled load;
  they do not upload, overwrite, or authorize a live CRM change.
\end{enumerate}

The package is versioned by the source-export hash and supporting
checksums. If an authorized correction changes a fact, a new snapshot is
built. The old snapshot remains available for audit; it is never edited
in place.

\subsection{What the files look like}\label{what-the-files-look-like}

The exact folder names may be tailored to the engagement, but the
contents have the following roles.

{\def\LTcaptype{none} % do not increment counter
\begin{longtable}[]{@{}
  >{\raggedright\arraybackslash}p{(\linewidth - 4\tabcolsep) * \real{0.3333}}
  >{\raggedright\arraybackslash}p{(\linewidth - 4\tabcolsep) * \real{0.3333}}
  >{\raggedright\arraybackslash}p{(\linewidth - 4\tabcolsep) * \real{0.3333}}@{}}
\toprule\noalign{}
\begin{minipage}[b]{\linewidth}\raggedright
Output
\end{minipage} & \begin{minipage}[b]{\linewidth}\raggedright
Typical form
\end{minipage} & \begin{minipage}[b]{\linewidth}\raggedright
What you use it for
\end{minipage} \\
\midrule\noalign{}
\endhead
\bottomrule\noalign{}
\endlastfoot
Final review package & JSON, CSV, XLSX, and HTML & Review open
questions, filter by priority or document, and return controlled
decisions. \\
Control results & JSON and human-readable summaries & Understand
arithmetic, validation, attribution, completeness, sampling, and
final-gate results. \\
Exception register & JSON and reviewer views & See every withheld
record, missing input, failed check, or unresolved value with its
reason. \\
Canonical export & Versioned JSON plus exclusion artifact & The
authoritative machine-readable approved-fact handoff. It is the source
for downstream loading and retrieval. \\
Load plan & Versioned, checksummed JSON & See the dependency order,
record counts, keys, and integrity checks required before loading
another system. \\
Retrieval snapshot & Immutable SQLite database & The local, read-only
source used by the MCP and REST API. It contains approved facts only. \\
CRM staging package & Ordered CSV files, API envelope, and manifest &
Give a target-system team a full, reconciled, no-send handoff across all
canonical tables. \\
Common CRM import package & Twelve ordered CSV files, mapping template,
source catalog, and write plan & Start target-CRM mapping and sandbox
testing using familiar business objects. \\
Optional source-intake package & Private original images, proposal
history, exceptions, manifest, and a derived PDF when complete & Hand
new source pages to the processing team. This is an evidence archive,
not an approved CRM import or a report. \\
\end{longtable}
}

\subsubsection{The Excel files}\label{the-excel-files}

Excel is used as a convenient review or handoff view, not as a place
where facts silently change.

\begin{itemize}
\tightlist
\item
  The \textbf{final review workbook} is for answering open questions.
  Its issued copy and returned copy are preserved separately. Returned
  decisions are proposals until an operator reviews the exact impact and
  reruns affected controls.
\item
  The \textbf{schema and CRM capability workbook} is an orientation aid.
  It explains tables, fields, API/MCP capabilities, sample report
  shapes, and import-file structure. It does not contain hidden
  authority to load a CRM or change an approved fact.
\item
  The \textbf{CRM import mapping template} is where a target-system
  owner completes vendor- and tenant-specific object and field mapping.
  Its included suggestions are starting points, not approved target
  mappings.
\end{itemize}

\subsection{What ``approved data''
means}\label{what-approved-data-means}

An approved row is not simply a row that an AI model read. Before it
reaches the canonical export or the retrieval snapshot, the system
requires applicable controls to be satisfied separately:

\begin{itemize}
\tightlist
\item
  source and page provenance are retained;
\item
  the record has a review-clear status;
\item
  applicable arithmetic is proved or explicitly not applicable;
\item
  the exhaustive final review queue has no open item for that record;
  and
\item
  every relevant attribution, completeness, and approval condition is
  recorded.
\end{itemize}

If a source, ledger export, payment export, mapping decision, or review
answer is unavailable, the package says so. It does not manufacture a
substitute fact. This is why a delivery can be valuable even when it is
not ready for a CRM load: it makes the outstanding business decisions
visible and bounded.

\subsection{What an MCP is}\label{what-an-mcp-is}

\textbf{MCP} stands for \textbf{Model Context Protocol}. In plain
language, it is a controlled connection that lets an AI assistant use
approved business tools. Instead of pasting spreadsheets into a chat,
the assistant can ask the approved retrieval snapshot specific, bounded
questions and show you the result.

Think of it as a librarian with a fixed catalog:

\begin{enumerate}
\def\labelenumi{\arabic{enumi}.}
\tightlist
\item
  You ask your assistant a business question in normal language.
\item
  The assistant selects an approved tool, such as search, exact lookup,
  sales analysis, or a standard report.
\item
  The tool reads one immutable, approved-fact snapshot.
\item
  The result includes source-export identity, checksums, and relevant
  source document IDs so it can be checked.
\end{enumerate}

The approved-data MCP does \textbf{not} let the assistant run arbitrary
SQL, browse raw AI responses, invent records, bypass the review gate, or
write to a CRM. Tool results do become part of the conversation context,
so the client must approve the selected AI provider's privacy,
retention, residency, and user-access settings before connecting real
data.

\subsection{How you can connect}\label{how-you-can-connect}

Choose the connection that fits how your team will work. The same
approved snapshot and semantic rules are used by every option; the
transport changes, not the facts or calculations.

{\def\LTcaptype{none} % do not increment counter
\begin{longtable}[]{@{}
  >{\raggedright\arraybackslash}p{(\linewidth - 6\tabcolsep) * \real{0.2500}}
  >{\raggedright\arraybackslash}p{(\linewidth - 6\tabcolsep) * \real{0.2500}}
  >{\raggedright\arraybackslash}p{(\linewidth - 6\tabcolsep) * \real{0.2500}}
  >{\raggedright\arraybackslash}p{(\linewidth - 6\tabcolsep) * \real{0.2500}}@{}}
\toprule\noalign{}
\begin{minipage}[b]{\linewidth}\raggedright
Your need
\end{minipage} & \begin{minipage}[b]{\linewidth}\raggedright
Connection
\end{minipage} & \begin{minipage}[b]{\linewidth}\raggedright
Suitable clients
\end{minipage} & \begin{minipage}[b]{\linewidth}\raggedright
What is required
\end{minipage} \\
\midrule\noalign{}
\endhead
\bottomrule\noalign{}
\endlastfoot
A quick, private desktop pilot & \textbf{Local stdio MCP} & Claude Code,
Claude Desktop, and compatible local desktop clients & The completed
run's local immutable snapshot and a private local registration. This is
the recommended first connection. \\
A controlled internal application or BI tool & \textbf{TLS/bearer REST
API} & An approved internal application or BI consumer & Explicit
enablement, TLS, dedicated rotated bearer secret, network controls,
monitoring, and client acceptance. This is an API, not an MCP
connector. \\
Web, iOS, or organization workspace use & \textbf{Remote Streamable HTTP
MCP} at a public \texttt{/\allowbreak{}mcp} URL & Claude
web/desktop/mobile after connector setup; selected ChatGPT workspace/API
surfaces & Public HTTPS, OAuth/OIDC authorization, exact
tenant/role/scope rules, approved workspace registration, monitoring,
and live acceptance. \\
\end{longtable}
}

\subsubsection{Local desktop connection: the simplest first
step}\label{local-desktop-connection-the-simplest-first-step}

After the engagement reaches the approved-data stage, the operator
registers a private local MCP configuration that points to the exact
local snapshot. In Claude Code, this is normally a local stdio
registration; Claude Desktop and other compatible desktop clients use
the same launcher through their local MCP configuration. The database
stays on the workstation, and the server is read-only.

This is the right starting point when one approved team is working from
the machine holding the completed snapshot. It does not make the service
public or available on a phone.

\subsubsection{Remote connection: needed for mobile and
web}\label{remote-connection-needed-for-mobile-and-web}

Claude mobile, claude.ai, and selected ChatGPT web, desktop, or mobile
experiences require a separately deployed remote MCP service. The
repository already supplies the remote approved-data server, but a real
client deployment still needs:

\begin{itemize}
\tightlist
\item
  one public HTTPS URL ending exactly in \texttt{/\allowbreak{}mcp};
\item
  a trusted TLS certificate and reverse proxy;
\item
  a real OAuth/OIDC authorization service with token introspection;
\item
  a specific tenant claim, approved roles, and
  \texttt{crm:read}/\texttt{crm:export} scopes;
\item
  approved logging, alerting, retention, backup, export cleanup, and
  incident response procedures; and
\item
  connector registration and representative acceptance tests in the
  selected Claude and/or ChatGPT workspace.
\end{itemize}

Mobile users normally use an already configured connector; they do not
create the deployment from the mobile app. Product-plan and workspace
availability can change, so the chosen provider surface is tested during
the deployment acceptance rather than assumed from protocol
compatibility.

\subsection{What you can ask for}\label{what-you-can-ask-for}

The system is designed for normal business questions. The assistant
chooses a bounded tool and returns a small, checkable result; it does
not send an entire large report into a chat conversation.

\subsubsection{Find a record when you do not know its
key}\label{find-a-record-when-you-do-not-know-its-key}

Use search for a company name, invoice number, address fragment,
acknowledgement number, job, project, or other source-visible phrase.

Examples:

\begin{itemize}
\tightlist
\item
  ``Find documents and records related to Acme Industrial.''
\item
  ``Search for invoice INV-1042 and show the approved source citation.''
\item
  ``Find the address containing 125 Harbor Drive.''
\item
  ``Which records mention ACK-7781 or job 44-19?''
\end{itemize}

The MCP uses \texttt{search\_\allowbreak{}crm\_\allowbreak{}records} for
this cross-table discovery. Once the right row is known,
\texttt{get\_\allowbreak{}crm\_\allowbreak{}record} retrieves the exact
canonical record by table and key.

\subsubsection{Pull an account, address, or contact
card}\label{pull-an-account-address-or-contact-card}

\texttt{get\_\allowbreak{}account\_\allowbreak{}card} is the preferred
way to ask for a business-ready view of a company or party. It joins
approved party names and roles with addresses, contacts, invoice/payment
activity, related source documents, and summary figures.

Examples:

\begin{itemize}
\tightlist
\item
  ``Show the contact card for Acme Industrial, including addresses and
  contacts.''
\item
  ``Give me the approved address and invoice history for customer
  CUST-104.''
\item
  ``Which source documents support this vendor's contact information?''
\end{itemize}

Money is separated by currency in account cards. The service will not
add US dollars, Canadian dollars, euros, or other currencies into one
misleading total.

\subsubsection{Slice and aggregate sales
safely}\label{slice-and-aggregate-sales-safely}

\texttt{analyze\_\allowbreak{}sales} is the governed analysis tool. It
can group up to three business dimensions, including company/customer,
vendor, region, selling location, month, quarter, year, currency,
acknowledgement number, and job number. It can also filter by date,
those same dimensions, and invoice-value bounds.

Examples:

\begin{itemize}
\tightlist
\item
  ``Show sales by company and region for Q2, separated by currency.''
\item
  ``Compare sales by selling location for this year.''
\item
  ``Break down customer sales for job J-204 by month.''
\item
  ``Which acknowledgement numbers contributed to Midwest region sales?''
\end{itemize}

Currency is always treated as a grouping dimension, even if it is not
mentioned in the question. That prevents unlike currencies from being
added together.

\subsubsection{Query one approved table}\label{query-one-approved-table}

\texttt{query\_\allowbreak{}crm\_\allowbreak{}records} is for a focused,
bounded slice of one table. It supports approved filters, selected
fields, stable sorting, and pagination. It does not execute formulas,
code, or arbitrary SQL.

Examples:

\begin{itemize}
\tightlist
\item
  ``List approved invoices for Acme Industrial with invoice number,
  date, amount due, and source document ID.''
\item
  ``Show the next 50 approved shipment records, using the available
  schema fields.''
\item
  ``Find addresses whose city starts with `New', then show the account
  cards for those parties.''
\end{itemize}

The last example needs more than one tool call: the one-table query does
not perform arbitrary joins from contacts to addresses. The assistant
must discover available fields and keys before choosing filters, not
invent columns.

For a larger reconciliation,
\texttt{export\_\allowbreak{}crm\_\allowbreak{}table} returns
checksummed pages rather than an unbounded conversation response.

\subsection{Standard reports available out of the
box}\label{standard-reports-available-out-of-the-box}

The retrieval layer supplies seven deterministic reports over approved
facts. Each report is bound to the snapshot hash and includes source
document IDs for its supporting data.

{\def\LTcaptype{none} % do not increment counter
\begin{longtable}[]{@{}
  >{\raggedright\arraybackslash}p{(\linewidth - 2\tabcolsep) * \real{0.5000}}
  >{\raggedright\arraybackslash}p{(\linewidth - 2\tabcolsep) * \real{0.5000}}@{}}
\toprule\noalign{}
\begin{minipage}[b]{\linewidth}\raggedright
Report
\end{minipage} & \begin{minipage}[b]{\linewidth}\raggedright
Typical use
\end{minipage} \\
\midrule\noalign{}
\endhead
\bottomrule\noalign{}
\endlastfoot
\texttt{account\_\allowbreak{}directory} & List approved
accounts/parties with key contact and address information. \\
\texttt{invoice\_\allowbreak{}register} & Review invoice-level amounts,
dates, parties, and source citations. \\
\texttt{receivables\_\allowbreak{}aging} & Review the current approved
balance snapshot by aging bucket as of a required date. It is not a
reconstructed historical ledger. \\
\texttt{sales\_\allowbreak{}by\_\allowbreak{}customer} & Compare
approved sales by customer, with currencies kept separate. \\
\texttt{sales\_\allowbreak{}by\_\allowbreak{}selling\_\allowbreak{}location}
& Compare approved sales by sales office, branch, or selling
location. \\
\texttt{attribution\_\allowbreak{}coverage} & See which approved invoice
values have a supported attribution and which do not. \\
\texttt{shipment\_\allowbreak{}performance} & Review approved shipment
activity and related operational measures available in the snapshot. \\
\end{longtable}
}

You can ask an assistant for these in ordinary language: ``Run the
invoice register for April,'' ``show sales by customer,'' or ``run
receivables aging as of June 30.'' The assistant maps the request to the
report and must preserve its parameters and snapshot identity in the
result.

\subsection{Downloads and large
exports}\label{downloads-and-large-exports}

Chat is useful for concise answers, not for sending thousands of rows
into a model conversation. The remote MCP therefore supports controlled
CSV and XLSX export jobs.

An authorized user can request one approved table or one standard report
in CSV or XLSX form. The server creates a bounded, short-lived download
with:

\begin{itemize}
\tightlist
\item
  the snapshot and source-export hash used at the start;
\item
  the requester identity in hashed form;
\item
  the exact table/report and filters requested;
\item
  row, column, and file checksums; and
\item
  an expiry time and owner-bound status lookup.
\end{itemize}

Spreadsheet-formula-like text is neutralized in exports. A different
user cannot inspect another user's job, a changed file fails
verification, and an expired link is refused. These exports create a
read-only report artifact; they never update canonical data or write to
a CRM.

\subsection{CRM inputs: what is ready and what still needs a
decision}\label{crm-inputs-what-is-ready-and-what-still-needs-a-decision}

\subsubsection{Adding photographs and proposed records
first}\label{adding-photographs-and-proposed-records-first}

The optional intake pilot can retain original PNG/JPEG page images and
let the connected, vision-capable assistant read them using the
pipeline's published field dictionary. It can preserve an invoice
reading or contact-card observation without changing the approved
business database. Unknown document families and unfamiliar labels stay
explicit, not forced into made-up fields.

\begin{enumerate}
\def\labelenumi{\arabic{enumi}.}
\tightlist
\item
  Agree which images and AI provider may receive the data; the operator
  enables intake separately from approved reporting.
\item
  A compatible application uploads the actual original bytes to a named
  session. Keep its receipt outside the chat. A chat attachment alone
  does not prove the service received it.
\item
  The assistant discovers the schema, reads every retained image, and
  proposes source-visible fields with original readings, page references
  and uncertainty.
\item
  The service checks the structure and returns \textbf{pending review}
  or \textbf{rejected}, retaining the original and every finding.
  Neither status is verification of the reading or permission to publish
  it.
\item
  Corrections append another proposal. An uncertain network retry uses
  the same request key; changed content uses a new one. A later proposal
  never erases an earlier failure or changes an approved record.
\end{enumerate}

For example: ``Read this session's retained invoice images. Propose the
invoice number, visible line items and amounts. Keep unreadable fields
and unfamiliar labels explicit, and show me the review receipt.'' For a
contact card, request source-visible observations, not an immediate
customer update.

The pilot accepts 1 to 20 pages per session: PNG/JPEG only, at most
10,000,000 original image bytes and 25,000,000 pixels per page. PDF,
HEIC, WebP, animations and image URLs are not direct upload types.
Preserve unsupported originals separately. Deployment request limits may
be smaller than the image limit. No mobile capture app, attachment
relay, session search, shared reviewer screen, or automatic record
publication is bundled. Actual upload and image reading must be tested
in the selected client; a working report connector proves neither.

\subsubsection{The source-only handoff}\label{the-source-only-handoff}

A local operator can export the session into a new private source
package. It retains every upload attempt, proposal version and rejection
with checksums. A complete usable image set also gets a session-named
PDF and an original-to-page map. The operator verifies the package
against a separately retained receipt, visually checks the PDF, and
starts the normal profiling/intake commands.

The PDF never replaces the originals. Display conversion keeps page
order and encoded orientation and shows transparency on white. It is not
a rescan, OCR result, or proof of original physical resolution. Missing
pages or conversion failures keep the archive blocked; a failed partial
PDF is never used. This is a controlled source archive, not
automatically a client-approved delivery.

Independent extraction, business controls, authorization and final
review still stand between these sources and a new approved snapshot.
The assistant's proposal is not automatically converted into a pipeline
vote and cannot change account cards, sales totals, reports or
CRM-import files. Full setup, fields, limits and recovery are in
\href{../references/visual-ingestion.md}{Visual Intake}.

\subsubsection{Approved-data CRM
preparation}\label{approved-data-crm-preparation}

The pipeline can prepare CRM-ready inputs without sending anything to a
CRM. This separation gives the target-system owner a chance to validate
mappings, permissions, duplicate behavior, automations, and rollback
before production.

\subsubsection{Full canonical staging
package}\label{full-canonical-staging-package}

The no-send staging package preserves all \textbf{28 canonical tables}
in dependency order. It includes ordered CSVs, a versioned API-operation
envelope, source and load-plan checksums, row counts, idempotency keys,
formula-safe display cells, and authoritative compact canonical JSON for
every row. It is designed for a receiver or integration team that needs
complete reconciliation across the full data model.

\subsubsection{Common CRM import
package}\label{common-crm-import-package}

For teams that work with familiar CRM objects, the optional common
package creates twelve ordered CSV files:

\begin{enumerate}
\def\labelenumi{\arabic{enumi}.}
\tightlist
\item
  accounts
\item
  account roles
\item
  addresses
\item
  contacts
\item
  products
\item
  sales locations
\item
  shipments
\item
  sales transactions
\item
  sales transaction lines
\item
  transaction charges
\item
  payments
\item
  payment applications
\end{enumerate}

It also includes a field-mapping template, a vendor-source catalog, a
write-readiness plan, exact row and package checksums, and complete
coverage accounting for all 28 canonical tables. A field not represented
by the common profile is explicitly marked as outside the profile rather
than silently dropped.

\subsubsection{What must happen before a live CRM
write}\label{what-must-happen-before-a-live-crm-write}

No live Salesforce, HubSpot, Dynamics/Dataverse, Zoho, or other CRM
writer is enabled by this package. A target-specific implementation
starts only after the client selects the vendor and tenant, supplies
current object/property metadata, approves mapping and credentials, and
proves the adapter in a sandbox.

The production-write acceptance must show stable external keys,
idempotent reruns, parent-before-child loading, explicit rejected-row
handling, count and financial reconciliation, relationship checks,
before-images/created IDs for rollback, and separate authorization for
the exact package hashes. Read-only MCP/API access does not grant
\texttt{crm:write} authority.

\subsection{A practical first-use
path}\label{a-practical-first-use-path}

For most teams, the safest and quickest sequence is:

\begin{enumerate}
\def\labelenumi{\arabic{enumi}.}
\tightlist
\item
  Review the final package and resolve or explicitly retain any open
  items.
\item
  Confirm the exact approved snapshot and its source-export checksum.
\item
  Start with a private local desktop MCP connection for representative
  questions and reports.
\item
  Validate account cards, invoice search, a sales slice, one standard
  report, and an exported CSV/XLSX against the delivered package.
\item
  If mobile, web, or multi-user access is needed, deploy the remote
  \texttt{/\allowbreak{}mcp} service with OAuth and complete
  workspace-specific acceptance.
\item
  If CRM loading is needed, complete target mapping and sandbox testing
  from the no-send package before authorizing a target-specific adapter.
\end{enumerate}

This order keeps a simple reporting pilot separate from a public
deployment or a live CRM change. It also gives business users an early
way to test whether the approved data answers their questions before the
organization commits to a larger integration.

\subsection{Questions this guide helps
answer}\label{questions-this-guide-helps-answer}

\begin{itemize}
\tightlist
\item
  \textbf{Can I find an address or contact quickly?} Yes, use search and
  the account card over the approved snapshot.
\item
  \textbf{Can I see sales by company, region, office, customer, job, or
  time period?} Yes, use governed sales analysis or the relevant
  standard report; currency is kept separate.
\item
  \textbf{Can I pull any approved record?} Yes, search when the key is
  unknown, then retrieve the exact canonical record by table and key.
\item
  \textbf{Can I download a report?} Yes, a remote deployment can create
  bounded, owner-bound CSV/XLSX jobs with checksums and expiry.
\item
  \textbf{Can an assistant change my CRM?} No.~Approved records remain
  read-only. Optional intake writes proposals to a separate journal, not
  canonical or target-CRM records. CRM inputs remain no-send until
  separate adapter implementation, authorization and acceptance.
\item
  \textbf{Can I submit a photo for a new record?} The optional pilot
  supports original PNG/JPEG upload and source-cited proposals through a
  compatible application. It stops at review, not publication.
\item
  \textbf{Can I ask for any possible analysis?} Shipped filters, account
  cards, sales dimensions and seven reports work within their limits.
  Custom joins, saved reports, forecasting and broader metrics need
  additional implementation and appropriate approved source data.
\item
  \textbf{Can the system tell me what it cannot prove?} Yes. Unresolved
  source, arithmetic, attribution, completeness, mapping, and review
  findings remain explicit rather than being hidden.
\end{itemize}

\subsection{Before connecting real
data}\label{before-connecting-real-data}

Use the following checklist with the engagement owner:

\begin{itemize}
\tightlist
\item[$\square$]
  The final review and applicable business gates are complete for the
  exact snapshot being served.
\item[$\square$]
  The team has chosen local desktop, internal REST, or remote MCP
  access.
\item[$\square$]
  The selected AI provider's privacy, retention, residency, and
  user-access terms are approved for the data returned in conversations.
\item[$\square$]
  Remote use has an approved domain, TLS, OAuth/OIDC, tenant/role/scope
  policy, monitoring, retention, and incident process.
\item[$\square$]
  The exact Claude and/or ChatGPT workspace connection has been tested
  with representative questions.
\item[$\square$]
  Any CRM load has a completed tenant mapping, sandbox evidence,
  reconciliation plan, rollback plan, and separate production
  authorization.
\item[$\square$]
  Optional intake has separately demonstrated byte upload, image
  reading, ownership checks, rejection/retry receipts, source export and
  pending-review wording; none is inferred from successful approved-data
  queries.
\end{itemize}

\subsection{Where to go next}\label{where-to-go-next}

\begin{itemize}
\tightlist
\item
  For the detailed output contract and approved retrieval rules, see
  \href{../references/canonical-deployment-retrieval.md}{Canonical
  Deployment, CRM Load, and Retrieval}.
\item
  For activation, connection, deployment, and acceptance steps, see
  \href{../references/mcp-production-integration.md}{MCP Production
  Integration}.
\item
  For the target-system import and write-readiness boundary, see
  \href{../references/crm-write-readiness.md}{CRM Write Readiness}.
\item
  For unimplemented capabilities, see the
  \href{../references/business-data-platform-roadmap.md}{vendor-neutral
  roadmap}.
\item
  For the broader evidence and review workflow, see the
  \href{CLIENT_OVERVIEW.md}{Client Overview} and
  \href{CLIENT_USER_GUIDE.md}{Client User Guide}.
\end{itemize}

The guiding rule is unchanged across every output: use the approved
snapshot for business answers, keep its citations and checksums with
important results, and treat a later correction as a new, controlled
version rather than an edit to history.

\end{document}
